\apendice{Roteiro de Entrevista (Exemplo)}
\label{ap:roteiro-de-entrevista-exemplo}

Este arquivo ilustra como fica um apêndice já preenchido -- neste caso, o roteiro de um instrumento de coleta de dados citado na Metodologia (\autoref{chap:metodologia}). Substitua pelo apêndice real do seu trabalho (questionário aplicado, código-fonte, tabelas de dados brutos, roteiro de entrevista etc.).

\begin{enumerate}
	\item Substitua pela primeira pergunta/etapa do seu instrumento de coleta.
	\item Substitua pela segunda pergunta/etapa do seu instrumento de coleta.
	\item Substitua pela terceira pergunta/etapa do seu instrumento de coleta.
\end{enumerate}

Se o seu trabalho não tiver necessidade de apêndices, apague este arquivo (e os demais desta pasta), removendo também as respectivas linhas \texttt{\textbackslash input\{...\}} e, se nenhum apêndice restar, a linha \texttt{\textbackslash imprimirapendices} do documento.tex.
